%%
%% preamble.tex — Shared LaTeX preamble for all papers in the folio.
%%
%% This file is a TEMPLATE.  Paper-specific content (title, authors,
%% macros, abstract, chapter \input lines) is generated by the content
%% pipeline from each paper's .ts manifest.
%%
%% Do NOT put paper-specific definitions here.
%%

\documentclass[11pt,a4paper]{report}

% ── Packages ──────────────────────────────────────────────────────────────────
\usepackage[utf8]{inputenc}
\usepackage[T1]{fontenc}
\usepackage{amsmath}
\usepackage{amssymb}
\usepackage{amsthm}
\usepackage{mathtools}
\usepackage{geometry}
\usepackage{graphicx}
\usepackage{cite}
\usepackage{ccicons}          % CC icon glyphs
\usepackage{titlesec}         % Compact heading styles
\usepackage{tikz}             % Inline figures
\usepackage{tikz-cd}          % Commutative diagrams (tikzcd environment)
\usepackage[dvipsnames]{xcolor}  % Extended colour names (teal, etc.)
\usepackage{booktabs}            % Professional table rules (\toprule, \midrule, \bottomrule)
\usepackage{enumitem}            % Customised list labels (label=, leftmargin=)
\setlist{noitemsep,topsep=0.5ex} % Condensed layout for lists
\usepackage{mathrsfs}            % \mathscr script letters
\usepackage{textcomp}            % \textdegree, \textbullet, etc.
\usepackage{newunicodechar}      % Map Unicode glyphs to LaTeX commands
\usepackage{algorithm}           % Float environment for pseudocode (algorithm)
\usepackage{algpseudocode}       % Algorithmic constructs (\State, \For, \If, …)
\usepackage[normalem]{ulem}      % latexdiff markup (\sout/\uwave); normalem keeps \emph italic
\usepackage{microtype}           % Typographic enhancements (character protrusion, font expansion)
\usepackage{adjustbox}           % Scale over-wide tables down to \textwidth (loaded after microtype)
\setlength{\emergencystretch}{3em} % Last-resort inter-word stretch for unbreakable long words/code/URLs
\usepackage{hyperref}            % Loaded last (after cite, xcolor, etc.)
\usepackage{xurl}                % MUST be after hyperref: break long URLs at any character
% KNOWN RESIDUAL: xurl/hyperref only break URLs typeset via \url{}/\href{} — e.g.
% author-note `see` links and the margin-annotation macros below. Bare URLs in
% running prose are emitted as PLAIN TEXT by the renderer (gfm-autolink-literal is
% disabled in content/pipeline/render-latex.ts to avoid OOM on dotted identifiers),
% so they are NOT broken here and remain a residual Overfull \hbox source. Folded in
% with the deeper residuals (wide display math, long \texttt identifiers) for a
% follow-up pass — to cover them, wrap bare prose URLs in \url{} at render time.

% ── Unicode character mappings ───────────────────────────────────────────────
% Allow Unicode symbols in .md content to pass through pdflatex.
\newcommand{\unimathsafe}[1]{\ifmmode #1\else $#1$\fi}
\newunicodechar{✓}{\unimathsafe{\checkmark}}
\newunicodechar{✗}{\unimathsafe{\times}}
\newunicodechar{✔}{\unimathsafe{\checkmark}}
\newunicodechar{★}{\unimathsafe{\bigstar}}
\newunicodechar{°}{\ifmmode ^{\circ}\else $^{\circ}$\fi}
\newunicodechar{ż}{\ifmmode\dot{z}\else\.z\fi}% Latin z-dot (citations/prose names)
\newunicodechar{×}{\unimathsafe{\times}}
\newunicodechar{←}{\unimathsafe{\leftarrow}}
\newunicodechar{→}{\unimathsafe{\rightarrow}}
\newunicodechar{≤}{\unimathsafe{\leq}}
\newunicodechar{≥}{\unimathsafe{\geq}}
\newunicodechar{≈}{\unimathsafe{\approx}}
\newunicodechar{∞}{\unimathsafe{\infty}}
\newunicodechar{Δ}{\unimathsafe{\Delta}}
\newunicodechar{α}{\unimathsafe{\alpha}}
\newunicodechar{β}{\unimathsafe{\beta}}
\newunicodechar{γ}{\unimathsafe{\gamma}}
\newunicodechar{δ}{\unimathsafe{\delta}}
\newunicodechar{ε}{\unimathsafe{\varepsilon}}
\newunicodechar{ζ}{\unimathsafe{\zeta}}
\newunicodechar{η}{\unimathsafe{\eta}}
\newunicodechar{θ}{\unimathsafe{\theta}}
\newunicodechar{ι}{\unimathsafe{\iota}}
\newunicodechar{κ}{\unimathsafe{\kappa}}
\newunicodechar{λ}{\unimathsafe{\lambda}}
\newunicodechar{μ}{\unimathsafe{\mu}}
\newunicodechar{ν}{\unimathsafe{\nu}}
\newunicodechar{ξ}{\unimathsafe{\xi}}
\newunicodechar{π}{\unimathsafe{\pi}}
\newunicodechar{ρ}{\unimathsafe{\rho}}
\newunicodechar{σ}{\unimathsafe{\sigma}}
\newunicodechar{τ}{\unimathsafe{\tau}}
\newunicodechar{υ}{\unimathsafe{\upsilon}}
\newunicodechar{φ}{\unimathsafe{\varphi}}
\newunicodechar{χ}{\unimathsafe{\chi}}
\newunicodechar{ψ}{\unimathsafe{\psi}}
\newunicodechar{ω}{\unimathsafe{\omega}}
\newunicodechar{Γ}{\unimathsafe{\Gamma}}
\newunicodechar{Λ}{\unimathsafe{\Lambda}}
\newunicodechar{Σ}{\unimathsafe{\Sigma}}
\newunicodechar{Π}{\unimathsafe{\Pi}}
\newunicodechar{Φ}{\unimathsafe{\Phi}}
\newunicodechar{Ψ}{\unimathsafe{\Psi}}
\newunicodechar{Ω}{\unimathsafe{\Omega}}
\newunicodechar{≠}{\unimathsafe{\neq}}
\newunicodechar{⊙}{\unimathsafe{\odot}}
\newunicodechar{●}{\textbullet}
\newunicodechar{○}{$\circ$}
\newunicodechar{…}{\ldots}
\newunicodechar{⚠}{\textbf{!}}
% Subscript digits (U+2080–U+2089)
\newunicodechar{₀}{\textsubscript{0}}
\newunicodechar{₁}{\textsubscript{1}}
\newunicodechar{₂}{\textsubscript{2}}
\newunicodechar{₃}{\textsubscript{3}}
\newunicodechar{₄}{\textsubscript{4}}
\newunicodechar{₅}{\textsubscript{5}}
\newunicodechar{₆}{\textsubscript{6}}
\newunicodechar{₇}{\textsubscript{7}}
\newunicodechar{₈}{\textsubscript{8}}
\newunicodechar{₉}{\textsubscript{9}}
% Superscript digits (U+2070, U+2074–U+2079)
\newunicodechar{⁰}{\textsuperscript{0}}
\newunicodechar{⁴}{\textsuperscript{4}}
\newunicodechar{⁵}{\textsuperscript{5}}
\newunicodechar{⁶}{\textsuperscript{6}}
\newunicodechar{⁷}{\textsuperscript{7}}
\newunicodechar{⁸}{\textsuperscript{8}}
\newunicodechar{⁹}{\textsuperscript{9}}
% Additional maths/arrow symbols (U+2194, U+2208, U+220F, U+2212, U+211D)
\newunicodechar{↔}{\unimathsafe{\leftrightarrow}}
\newunicodechar{∈}{\unimathsafe{\in}}
\newunicodechar{∏}{\unimathsafe{\prod}}
\newunicodechar{−}{\unimathsafe{-}}
\newunicodechar{ℝ}{\unimathsafe{\mathbb{R}}}
% Symbols appearing in prose/inline-code that lack T1 glyphs (U+2113, U+2261)
\newunicodechar{ℓ}{\unimathsafe{\ell}}
\newunicodechar{≡}{\unimathsafe{\equiv}}
% Additional math symbols and emoji that appear in prose
\newunicodechar{ᵢ}{\textsubscript{i}}            % U+1D62
\newunicodechar{′}{\unimathsafe{{}^{\prime}}}    % U+2032 prime
\newunicodechar{⁻}{\textsuperscript{$-$}}        % U+207B superscript minus
\newunicodechar{ℂ}{\unimathsafe{\mathbb{C}}}     % U+2102
\newunicodechar{ℕ}{\unimathsafe{\mathbb{N}}}     % U+2115
\newunicodechar{ℚ}{\unimathsafe{\mathbb{Q}}}     % U+211A
\newunicodechar{ℤ}{\unimathsafe{\mathbb{Z}}}     % U+2124
\newunicodechar{⇒}{\unimathsafe{\Rightarrow}}    % U+21D2
\newunicodechar{∀}{\unimathsafe{\forall}}        % U+2200
\newunicodechar{∃}{\unimathsafe{\exists}}        % U+2203
\newunicodechar{∎}{\unimathsafe{\blacksquare}}   % U+220E QED
\newunicodechar{∑}{\unimathsafe{\sum}}           % U+2211
\newunicodechar{∘}{\unimathsafe{\circ}}          % U+2218
\newunicodechar{∧}{\unimathsafe{\wedge}}         % U+2227
\newunicodechar{∨}{\unimathsafe{\vee}}           % U+2228
\newunicodechar{≅}{\unimathsafe{\cong}}          % U+2245
\newunicodechar{⊂}{\unimathsafe{\subset}}        % U+2282
\newunicodechar{⊆}{\unimathsafe{\subseteq}}      % U+2286
\newunicodechar{⊕}{\unimathsafe{\oplus}}         % U+2295
\newunicodechar{⊗}{\unimathsafe{\otimes}}        % U+2297
\newunicodechar{⟺}{\unimathsafe{\Longleftrightarrow}} % U+27FA
\newunicodechar{⧸}{\unimathsafe{/}}              % U+29F8 big solidus
% Status emoji used in tables/prose: render as text labels
\newunicodechar{⏳}{\textbf{[pending]}}            % U+23F3 hourglass
\newunicodechar{✅}{\textbf{[done]}}               % U+2705 check mark
\newunicodechar{🔓}{\textbf{[open]}}               % U+1F513 unlocked
\newunicodechar{🟡}{\textbf{[!]}}                  % U+1F7E1 yellow circle
% Additional status/math Unicode added to fix LaTeX errors in generated chapters
\newunicodechar{⚪}{\textbf{[--]}}                 % U+26AA white circle (pending)
\newunicodechar{🟢}{\textbf{[OK]}}                 % U+1F7E2 green circle
\newunicodechar{🔴}{\textbf{[X]}}                  % U+1F534 red circle
\newunicodechar{🟠}{\textbf{[~]}}                  % U+1F7E0 orange circle
\newunicodechar{⟵}{\unimathsafe{\longleftarrow}}  % U+27F5 long left arrow
\newunicodechar{⟹}{\unimathsafe{\Longrightarrow}} % U+27F9 long double right arrow
\newunicodechar{↦}{\unimathsafe{\mapsto}}          % U+21A6 maps to
\newunicodechar{─}{\texttt{-}}                     % U+2500 box drawing horizontal
\newunicodechar{│}{\texttt{|}}                     % U+2502 box drawing vertical
\newunicodechar{⋆}{\unimathsafe{\star}}            % U+22C6 star operator
\newunicodechar{□}{\unimathsafe{\square}}          % U+25A1 white square
\newunicodechar{◄}{\unimathsafe{\triangleleft}}    % U+25C4 left-pointing triangle
\newunicodechar{▼}{\unimathsafe{\triangledown}}    % U+25BC down-pointing triangle
\newunicodechar{►}{\unimathsafe{\triangleright}}   % U+25BA right-pointing triangle
\newunicodechar{▲}{\unimathsafe{\triangle}}        % U+25B2 up-pointing triangle
\newunicodechar{ℏ}{\unimathsafe{\hbar}}            % U+210F h-bar
\newunicodechar{√}{\unimathsafe{\surd}}            % U+221A square root
\newunicodechar{̄}{}                                % U+0304 combining macron (no-op; use \bar in math)
\newunicodechar{̂}{}                                % U+0302 combining circumflex (no-op; use \hat in math)
\newunicodechar{·}{\unimathsafe{\cdot}}            % U+00B7 middle dot
\newunicodechar{†}{\unimathsafe{\dagger}}          % U+2020 dagger
% Math symbols found in generated chapters (2026-05-29 PDF-compile fix)
\newunicodechar{⊢}{\unimathsafe{\vdash}}           % U+22A2 turnstile (λ ⊢ n)
\newunicodechar{↓}{\unimathsafe{\downarrow}}       % U+2193 downward arrow
\newunicodechar{⟨}{\unimathsafe{\langle}}          % U+27E8 math left angle
\newunicodechar{⟩}{\unimathsafe{\rangle}}          % U+27E9 math right angle
\newunicodechar{•}{\unimathsafe{\cdot}}            % U+2022 bullet (smul)
\newunicodechar{𝓑}{\unimathsafe{\mathcal{B}}}      % U+1D4D1 script B (braid word)

% ── Custom symbols ───────────────────────────────────────────────────────────
\newcommand{\bigbowtie}{\mathop{\triangleright\!\!\!\triangleleft}}
% Contact form (pseudo-Hermitian structure) in the QOU archimedean 4-tuple.
% \varTheta is the upright capital theta, visually distinct from the
% lowercase contact-form letter θ and from italic \Theta.
% Some recent TeX Live distributions predefine \varTheta; redefine safely.
\providecommand{\varTheta}{}
\renewcommand{\varTheta}{\mathrm{\Theta}}
% Reeb vector field variants used in contact/brane geometry
\newcommand{\xid}{{\xi_{d}}}   % ξ_d — Reeb field (d-component)
\newcommand{\xiZ}{{\xi_{Z}}}   % ξ_Z — Reeb field (Z-component)
\newcommand{\xiA}{{\xi_{A}}}   % ξ_A — Reeb field (A-component)
% ── Stubs for compound macros (missing spaces in source content) ──────────────
% These arise from LaTeX commands written without a separating space/brace.
\newcommand{\pp}{p}             % prime p — used in p-adic content (\mathbb{Q}_\pp, \mathbb{F}_\pp, v_\pp). Italic per math-mode convention; not \mathfrak{p} which is reserved for the proton (CLAUDE.md §7aa).
\newcommand{\pid}{\pi d}
\newcommand{\nuk}{\nu_{k}}
\newcommand{\varepsilonk}{\varepsilon_{k}}
\newcommand{\varepsilonq}{\varepsilon_{q}}
\newcommand{\foralln}{\forall n}
\newcommand{\omegadV}{\omega\,dV}
\newcommand{\omegads}{\omega\,ds}
\newcommand{\Gammat}{\Gamma t}
\newcommand{\psiK}{\psi_{K}}
\newcommand{\pic}{\pi}
\newcommand{\pia}{\pi a}
\newcommand{\lambdaL}{\lambda L}
\newcommand{\inftydt}{\infty\,dt}
% \unicode{codepoint} — ConTeXt-style; used in appendix-surreals for ☉ (U+2609)
\newcommand{\unicode}[1]{\ensuremath{\odot}}
% Young-diagram macro stub (ytableau not loaded; renders partition as list)
\newcommand{\yng}[1]{{\scriptsize[#1]}}

% Unicode maths symbols that appear in plain text (outside $ … $)
% are handled by \newunicodechar in the block above.

\geometry{
  left=0.8in, right=1.5in, top=0.8in, bottom=0.8in,
  marginparwidth=1.1in, marginparsep=0.15in,
}
\usepackage{marginnote}
\renewcommand*{\marginfont}{\scriptsize}

% ── Suppress hyperref math-in-bookmark warnings ─────────────────────────────
% Math in theorem titles generates "Token not allowed in PDF string" warnings.
% These are harmless (hyperref strips the math for bookmarks) but noisy.
\hypersetup{pdfencoding=auto, psdextra}
\pdfstringdefDisableCommands{%
  \def\({}%
  \def\){}%
  \let\mathbf\relax
  \let\mathrm\relax
  \let\mathcal\relax
  \let\mathfrak\relax
  \let\mathbb\relax
  \let\mathscr\relax
  \let\boldsymbol\relax
  \let\textstyle\relax
  \let\Sigma\relax
  \let\Delta\relax
  \let\not\relax
  \let\to\relax
  \let\bigbowtie\relax
}
% Silence the "Token not allowed" warnings entirely — they are
% informational and produce correct bookmarks regardless.
\makeatletter
\renewcommand{\HyPsd@Warning}[1]{}
\makeatother

% ── Book-style chapter / section formatting ────────────────────────────────
% Chapter page breaks are handled by the renderer (it prepends \clearpage
% to every \input chapter — see content/pipeline/render-latex.ts
% renderChapter). The chapter numeral here is rendered very large,
% left-aligned, with whitespace to the right — book convention. The
% chapter title sits below the numeral. Sections and subsections stay
% compact with normal-size fonts.
\titleclass{\chapter}{straight}       % renderer-driven page breaks
\titleformat{\chapter}[display]%
  {\normalfont\bfseries\raggedright}%
  {%
    \fontsize{96pt}{96pt}\selectfont\bfseries\thechapter%
    \hfill%
  }%
  {2em}%
  {\Huge\bfseries\raggedright}%
\titlespacing*{\chapter}{0pt}{0pt}{2em}

% Unnumbered chapters (introduction, glossary, appendices) — same display
% block but no large numeral, since \chapter* has no number.
\titleformat{name=\chapter,numberless}[block]%
  {\normalfont\Huge\bfseries\raggedright}%
  {}%
  {0pt}%
  {\Huge\bfseries\raggedright}%
\titlespacing*{name=\chapter,numberless}{0pt}{0pt}{2em}

\titleformat{\section}[hang]{\normalsize\bfseries}{%
  \thesection}{0.5em}{}
\titlespacing*{\section}{0pt}{1.0ex plus 0.3ex minus 0.2ex}{0.5ex plus 0.1ex}

\titleformat{\subsection}[hang]{\normalsize\bfseries}{%
  \thesubsection}{0.5em}{}
\titlespacing*{\subsection}{0pt}{0.8ex plus 0.3ex minus 0.1ex}{0.4ex plus 0.1ex}

% Show only chapters and sections in the table of contents — no subsections.
\setcounter{tocdepth}{1}

% ── Theorem environments ───────────────────────────────────────────────────────
% Condensed spacing to reduce vertical footprint
\newtheoremstyle{condensed}
  {0.8ex plus 0.2ex minus 0.1ex} % Space above
  {0.8ex plus 0.2ex minus 0.1ex} % Space below
  {\itshape}                     % Body font (italic for theorems)
  {}                             % Indent amount
  {\bfseries}                    % Theorem head font
  {.}                            % Punctuation after theorem head
  {0.5em}                        % Space after theorem head
  {}                             % Theorem head spec

\newtheoremstyle{condenseddef}
  {0.8ex plus 0.2ex minus 0.1ex} % Space above
  {0.8ex plus 0.2ex minus 0.1ex} % Space below
  {\normalfont}                  % Body font (roman for definitions)
  {}                             % Indent amount
  {\bfseries}                    % Theorem head font
  {.}                            % Punctuation after theorem head
  {0.5em}                        % Space after theorem head
  {}                             % Theorem head spec

\theoremstyle{condensed}
\newtheorem{theorem}{Theorem}[section]
\newtheorem{lemma}[theorem]{Lemma}
\newtheorem{proposition}[theorem]{Proposition}
\newtheorem{corollary}[theorem]{Corollary}
\newtheorem{conjecture}[theorem]{Conjecture}

\theoremstyle{condenseddef}
\newtheorem{definition}[theorem]{Definition}
\newtheorem{example}[theorem]{Example}
\newtheorem{remark}[theorem]{Remark}
\newtheorem{observation}[theorem]{Observation}

% `algorithmblock` (theorem-like) — separate from the `algorithm`
% float package, which is reserved for pseudocode floats elsewhere
% in the paper.
\newtheorem{algorithmblock}[theorem]{Algorithm}

% ── Proof environment formatting ───────────────────────────────────────────────
% Hollow square for QED, italic 'Proof' heading to distinguish from bold theorems
\renewcommand{\qedsymbol}{$\square$}
\makeatletter
\renewenvironment{proof}[1][\proofname]{\par
  \pushQED{\qed}%
  \normalfont \topsep0.8ex\@plus0.2ex\@minus0.1ex\relax
  \trivlist
  \item[\hskip\labelsep
        \itshape
    #1\@addpunct{.}]\ignorespaces
}{%
  \popQED\endtrivlist\@endpefalse
}
\makeatother

% ── LeanBlueprint macros ──────────────────────────────────────────────────────
% These macros annotate theorem-like environments with Lean 4 declaration
% linkages.  They are consumed by the proof-objects extraction pipeline
% and (optionally) by leanblueprint/plasTeX for dependency graph generation.
%
% Usage:
%   \begin{theorem}\label{thm:my-result}
%     \lean{QOU.Chapter.my_result}      % Lean declaration name
%     \leanok                            % proof is sorry-free
%     \uses{def:prerequisite}            % dependency labels
%     Statement here.
%   \end{theorem}
% Base URL for Lean documentation (doc-gen4 pages on GitHub Pages)
\newcommand{\leandocsbase}{https://litlfred.github.io/qou/lean/}
% GitHub source URL for Lean files
\newcommand{\leansrcbase}{https://github.com/litlfred/qou/blob/main/lean/}
% Repo-scoped GitHub code-search base.  We only carry the declaration name
% (not a reliable file path), so the ∀ status mark links here: it lands on a
% GitHub view of the declaration's source within the repo.
\newcommand{\leansrcsearchbase}{https://github.com/litlfred/qou/search?q=}
%
% ── Lean formalisation-status palette ───────────────────────────────────────
% Single source of truth for the three-state ∀ colour code.  Tweak here.
% Orange is deliberately avoided: it reads poorly on white and sits too close
% to red.  The short text tag (stub/draft/ok) carries the status regardless of
% how distinguishable the colours are on a given display/print.
\colorlet{leanstubcolor}{BrickRed}     % stubbed  — sorry/placeholder, or QA-vacuous/trivial
\colorlet{leandraftcolor}{Purple}      % drafted  — genuine statement in Lean, not yet sorry-free
\colorlet{leanokcolor}{ForestGreen}    % compiled — built sorry-free
%
% Internal flag: tracks whether \lean was called in the current environment.
\newif\ifleantagged
%
% \lean{Decl.Name} — renders a clickable, bold ∀ icon linking to the Lean
% declaration's source on GitHub.  Standalone, this is the "drafted" state
% (statement stated in Lean); a following \leanok adds the green compiled mark.
\newcommand{\lean}[1]{%
  \global\leantaggedtrue%
  \texorpdfstring{%
    \href{\leansrcsearchbase#1}{%
      \textsuperscript{\textcolor{leandraftcolor}{$\boldsymbol{\forall}$}}%
    }%
  }{(Lean: #1)}%
}
%
% \leanok — marks environment as fully formalized (sorry-free).
% Renders a green checkmark.
\newcommand{\leanok}{%
  \texorpdfstring{%
    \textsuperscript{\textcolor{ForestGreen}{\checkmark}}%
  }{[proved]}%
}
%
% \notready — marks environment as not ready to formalize.
% Renders a red cross.
\newcommand{\notready}{%
  \texorpdfstring{%
    \textsuperscript{\textcolor{BrickRed}{$\times$}}%
  }{[not ready]}%
}
%
% \mathlibok — marks declarations merged into mathlib.
% Renders a dark-green double checkmark.
\newcommand{\mathlibok}{%
  \texorpdfstring{%
    \textsuperscript{\textcolor{OliveGreen}{\checkmark\kern-0.3em\checkmark}}%
  }{[mathlib]}%
}
\newcommand{\uses}[1]{}            % \uses{label1, label2, ...}
\newcommand{\proves}[1]{}          % \proves{thm:label} (in proof environments)
%
% ── Margin annotations: source links + issue tracker + structured comments ──
%
% Every labelled content block, section, and chapter carries a compact
% strip of clickable icons in the right margin:
%
%   ∀  → Lean docs page for this block's declaration (if any)
%   ∇  → .md source on GitHub (the narrative form, authoritative per CLAUDE.md)
%   #  → GitHub issues filtered by the block label (existing discussion)
%   +  → "Add comment / open new issue" pre-filled with the block label
%
% The macros below are pure LaTeX, no font dependencies. They are
% emitted by content/pipeline/render-latex.ts after each \label{}.
% Empty arguments hide the corresponding icon, so a block without a
% Lean ref still gets ∇ # +, etc.

% Base URLs.  Override at compile time via \renewcommand{...}.
\newcommand{\trackerbase}{https://github.com/litlfred/qou/issues?q=}
\newcommand{\sourcebase}{https://github.com/litlfred/qou/blob/main/content/}
\newcommand{\commentbase}{https://github.com/litlfred/qou/issues/new?title=Comment+on+}

% \blockannot{label}{leanDecl}{mdRelPath}
% \blockannotok{label}{leanDecl}{mdRelPath}   — same but adds green ✓
%                                                indicating the Lean proof
%                                                is sorry-free.
% Compact horizontal strip of icons in the right margin. Any of the three
% positional args may be empty — that icon is simply omitted. The leading
% \mbox{} anchors the marginnote to a non-floating reference point in the
% body line so adjacent margin notes don't collide.
%
% NOTE: \makeatletter is required here so that \blockannot@i (and the
% internal \@tmp / \@empty refs in its body) are tokenised with @ as a
% letter (catcode 11).  Without it, \blockannot@i is tokenised as
% \blockannot (catcode-12 @ stops the control-word scan), which causes
% "Command \blockannot already defined" and a subsequent parameter-stack
% overflow when \blockannot recurses into itself.
% \leanstatusmark{status}{decl} — the Lean formalisation indicator: a BOLD,
% colour-coded ∀ plus a short text tag, the whole thing hyperlinked to the
% declaration's source on GitHub.  The tag makes the status legible even when
% the colours are hard to tell apart on a given display or in print.
%   status = stubbed  → red    ∀ stub   (sorry/placeholder, or vacuous/trivial per QA)
%   status = drafted  → purple ∀ draft  (genuine statement in Lean, not yet sorry-free)
%   status = compiled → green  ∀ ok     (built sorry-free)
\makeatletter
\newcommand{\leanstatusmark}[2]{%
  \def\lsm@s{#1}%
  \def\lsm@stub{stubbed}\def\lsm@comp{compiled}%
  \ifx\lsm@s\lsm@stub
    \def\lsm@col{leanstubcolor}\def\lsm@tag{stub}%
  \else\ifx\lsm@s\lsm@comp
    \def\lsm@col{leanokcolor}\def\lsm@tag{ok}%
  \else
    \def\lsm@col{leandraftcolor}\def\lsm@tag{draft}%
  \fi\fi
  \href{\leansrcsearchbase #2}{%
    \textcolor{\lsm@col}{$\boldsymbol{\forall}$\,\textbf{\tiny\lsm@tag}}%
  }%
}
% \blockannot[status]{label}{leanDecl}{mdRelPath}
%   Optional status keyword (default "drafted") drives the ∀ colour + tag.
% \blockannotok{label}{leanDecl}{mdRelPath}  ==  \blockannot[compiled]{...}
%   (kept for back-compat with previously generated content).
\newcommand{\blockannot}[4][drafted]{\blockannot@i{#1}{#2}{#3}{#4}}
\newcommand{\blockannotok}[3]{\blockannot@i{compiled}{#1}{#2}{#3}}
\newcommand{\blockannot@i}[4]{%
  \mbox{}\marginnote{%
    \raggedright\scriptsize
    \ifx\relax#3\relax\else
      \leanstatusmark{#1}{#3}\,%
    \fi
    \ifx\relax#4\relax\else
      \href{\sourcebase #4}{\textcolor{NavyBlue}{$\nabla$}}\,%
    \fi
    \ifx\relax#2\relax\else
      \href{\trackerbase #2}{\textcolor{NavyBlue}{\#}}\,%
      \href{\commentbase #2}{\textcolor{NavyBlue}{$+$}}%
    \fi
  }%
}
\makeatother

% \chapterannot{chapterDir}{trackerLabel}
% Chapter-level annotation: link to the chapter directory + tracker.
\newcommand{\chapterannot}[2]{%
  \mbox{}\marginnote{%
    \raggedright\scriptsize
    \href{\sourcebase #1/}{\textcolor{NavyBlue}{$\nabla$ src}}\\[-0.3ex]%
    \href{\trackerbase #2}{\textcolor{NavyBlue}{\# issues}}\\[-0.3ex]%
    \href{\commentbase #2}{\textcolor{NavyBlue}{$+$ note}}%
  }%
}

% \sectionannot{label} — section-level annotation: tracker only.
\newcommand{\sectionannot}[1]{%
  \marginnote{%
    \raggedright\scriptsize
    \href{\trackerbase #1}{\textcolor{NavyBlue}{\#}}\,%
    \href{\commentbase #1}{\textcolor{NavyBlue}{$+$}}%
  }%
}

% Back-compat: \tracker{label} still works (subsumed by \blockannot but kept
% so old hand-written .md content references don't break).
\newcommand{\tracker}[1]{\blockannot{#1}{}{}}

% \bcomment[author]{text} — structured human/agent comment in the margin.
% `text` empty → renders an "(add note)" stub. Optional [author] tag appears
% as a small "— Author" suffix.  Renders in BrickRed to distinguish from
% the auto-generated annotation icons.
\newcommand{\bcomment}[2][]{%
  \marginnote{%
    \raggedright\scriptsize\textcolor{BrickRed}{%
      \textbf{note:}~%
      \ifx\relax#2\relax\textit{(add note)}\else#2\fi%
      \ifx\relax#1\relax\else~\textemdash~\textit{#1}\fi%
    }%
  }%
}
%
% ── Glossary terms: \defterm{} / \refterm{} ─────────────────────────────────
% PDF rendering is intentionally clean — the HTML viewer carries the
% interactive affordance (dotted underline on hover, click-to-jump).
%
%   \defterm{slug}{Visible}  defining occurrence: bold-italic + invisible
%                            \phantomsection\label{term:slug} target
%   \refterm{slug}{Visible}  referencing occurrence: invisible
%                            \hyperref[term:slug]{Visible} (no decoration)
%
% Single-arg forms \defterm{slug} / \refterm{slug} use the slug as the
% visible label (the common case where slug == display text).
% NOTE: defined as \DeclareRobustCommand (not \newcommand) so the
% macros do not expand prematurely in moving arguments such as
% section titles.  The bare \newcommand form caused titlesec to
% recurse expanding `\refterm{descent}` inside a subsubsection
% title until the input stack overflowed — Build & Publish #3262
% failed with "TeX capacity exceeded [input stack size=10000]"
% expanding that exact macro on l.3752 of the rendered chapter.
% \texorpdfstring keeps PDF bookmarks safe by stripping the link.
\makeatletter
\DeclareRobustCommand{\defterm}[1]{\@ifnextchar\bgroup{\defterm@two{#1}}{\defterm@one{#1}}}
\newcommand{\defterm@one}[1]{\phantomsection\label{term:#1}\textbf{\emph{#1}}}
\newcommand{\defterm@two}[2]{\phantomsection\label{term:#1}\textbf{\emph{#2}}}
\DeclareRobustCommand{\refterm}[1]{\@ifnextchar\bgroup{\refterm@two{#1}}{\refterm@one{#1}}}
\newcommand{\refterm@one}[1]{\texorpdfstring{\hyperref[term:#1]{#1}}{#1}}
\newcommand{\refterm@two}[2]{\texorpdfstring{\hyperref[term:#1]{#2}}{#2}}
\makeatother
%
% ── Proof-status legend (printed once in the document) ───────────────────────
\newcommand{\proofstatuslegend}{%
  \bigskip\noindent\textbf{Lean formalisation key} — the bold ∀ beside each
  block links to its source on GitHub; its colour and tag give the status:%
  \par\smallskip\noindent
  \begin{tabular}{@{}cl@{\qquad}cl@{}}
    \textcolor{leanstubcolor}{$\boldsymbol{\forall}$\,\textbf{\footnotesize stub}}
      & Stub: \texttt{sorry}/placeholder, or vacuous/trivial per QA
    & \textcolor{leandraftcolor}{$\boldsymbol{\forall}$\,\textbf{\footnotesize draft}}
      & Drafted: genuine statement in Lean, not yet sorry-free\\
    \textcolor{leanokcolor}{$\boldsymbol{\forall}$\,\textbf{\footnotesize ok}}
      & Compiled: built sorry-free
    & \textcolor{OliveGreen}{\checkmark\kern-0.3em\checkmark} & Merged into Mathlib\\
  \end{tabular}\bigskip
}

% ── Print mode: affiliations ──────────────────────────────────────────────────
% In compact mode (default), affiliations are omitted.
% Set \showaffiliationstrue for formal print view.
% The MCP render tool writes print-mode.tex to set this flag automatically.
\newif\ifshowaffiliations
% The per-render print-mode flag (the read of print-mode.tex) is emitted by
% content/pipeline/generate-main-tex.ts into the per-document section, NOT here.
% Keeping only \newif here (shared, static) leaves this preamble fully dumpable
% into a precompiled .fmt format; the flag must be set live per render, not
% frozen at format-dump time.
